(1) 体内不可压 $\nabla\cdot\boldsymbol v=0$ 且 $\boldsymbol v=\nabla\Phi$，故 $\nabla^2\Phi=0$。固底没有垂直速度，$\partial_z\Phi(x,-h,t)=0$。设 $\Phi(x,z,t)=F(z)e^{i(kx-\omega t)}$，空间方程 $F''-k^2F=0$。由底部条件可取
\[
F(z)=A\frac{\cosh[k(z+h)]}{\cosh(kh)}.
\]
在自由表面 $z=0$，$F(0)=A,\ F'(0)=Ak\tanh(kh)$。

(2) 小扰动自由表面的运动学条件是
\[
\partial_t\eta=\partial_z\Phi|_{z=0},
\]
动力学条件来自线性化 Bernoulli 方程并取表面压强恒定：
\[
\partial_t\Phi|_{z=0}+g\eta=0.
\]
取 $\eta=\eta_0 e^{i(kx-\omega t)}$，动力学条件给 $\eta_0=i\omega A/g$。代回运动学条件，
\[
-i\omega(i\omega A/g)=Ak\tanh(kh),
\]
即
\[
\boxed{\omega^2=gk\tanh(kh)}.
\]
只有实正根对应向右传播的正频波，流体密度因无表面张力且重力为唯一回复机制而不出现在色散式中。

(3) 相速度 $c_p=\omega/k=\sqrt{(g/k)\tanh(kh)}$。对 $\omega^2$ 求 $k$ 导数，
\[
2\omega\,\frac{d\omega}{dk}
=g[\tanh(kh)+kh\,\operatorname{sech}^2(kh)].
\]
因此群速度
\[
c_g=\frac{d\omega}{dk}
=\frac{c_p}{2}\left[1+\frac{2kh}{\sinh(2kh)}\right].
\]
对 $kh>0$，有 $0<2kh/\sinh(2kh)<1$，故 $c_p/2<c_g<c_p$，反映自由表面重力波的正常频散。

(4) 浅水 $kh\ll1$ 时 $\tanh(kh)\sim kh$，故 $\omega\simeq k\sqrt{gh}$、$c_p\simeq c_g\simeq\sqrt{gh}$，波速几乎与波数无关，波包近似无频散。深水 $kh\gg1$ 时 $\tanh(kh)\to1$，故 $\omega\simeq\sqrt{gk}$、$c_p\simeq\sqrt{g/k}$、$c_g\simeq c_p/2$，波包移动速度是相位速度的一半。若加入表面张力，短波色散将改变，不能继续把高频结果当作普适定律。
